% ============================================================
% 分片：s03b-singularity  对应原文 PDF 第 20--28 页（印刷页码 20--28）
% 负责章节：3.2 The "Singularity"
% 译者：t05-s32   审校：t03-s2b（逐式核对完成）   审校状态：二审
% ============================================================
% 范围说明：本分片自 PDF 第 20 页第一个词起（承接 s03a 末句 "suffice it to say"），
% 包含上一小节末尾的 "The Moral Of The Story" 小标题及其段落，随后为 3.2 全节，
% 至 PDF 第 28 页末为止；不取第 29 页任何内容。
% 原文脚注（第 27 页 ⋆、第 28 页 †）以 \footnote{} 保留，编号由 LaTeX 重排。
% 原文用直立体 d 表示微分，故微分记号统一写作 \mathrm{d}。
% ==== 原文 PDF 第 20 页（印刷页 20）====
\label{op:20}
% [跨页接续: 原文第19页末 "suffice it to say" 与第20页首 "here that" 属同一句；s03a 已译出 "这里只需指出"，故本句不再重译 "here"，直接接 "在 R 不变的情形下"] 
在 $R$ 不变的情形下，对于足够大的 $|r|$，积掉有质量场只会修正有效低能理论中的参数，而不会使系统的性质发生定性变化。

\subsection*{故事的寓意}

上述构造提示：与其说 Landau--Ginzburg 与 Calabi--Yau 是“等价的”，不如说它们是同一系统的两个不同的相。事实上，从这一观点看，Landau--Ginzburg 看起来像是 Calabi--Yau 向负的凯勒类的解析延拓。在作若干进一步分析之后，我们将重新审视这一故事的寓意。

\subsection{“奇点”}

为了从 Landau--Ginzburg 与 Calabi--Yau 模型的关系中得出任何精确的结论，我们必须处理在 $r=0$ 处把二者分隔开的那个表观奇点。

问题并不在于 $r=0$ 处可能存在相变。在我们的应用中，我们关心的或是在圆周上（紧致、有限体积）量子化这一理论，或是在紧致黎曼面上作路径积分。无论哪种方式，我们都看不到通常与相变相联系的那种奇点；后者本质上有赖于无限的世界面体积。我们可能看到的唯一奇点，将来自靶空间有效紧致性（effective compactness）的失效。

% [Tutor Note: 这里把“相变奇点”换成了“有效紧致性”这一判据。在紧致世界面或圆周上，配分函数是参数的解析函数，除非某个场方向的次能级集 $\{V<T\}$ 非紧、使谱出现连续部分，路径积分在无穷远处的区域不再被指数压低。热力学意义上的相变要求世界面体积趋于无穷（配分函数的非解析性只在热力学极限中出现），故有限体积下本就不会看到它；本节此后所说的“奇点”一律指有效紧致性失效，而非相变。]

例如，若 $\sigma$ 场不存在（我们将在 §6 中看到，可以在保持 $(0,2)$ 超对称的前提下将其消去），则在有限世界面体积下，$r=0$ 处不可能出现奇点。因为若没有 $\sigma$ 场，$r=0$ 时经典玻色势唯一可能为零的点就是 $s_i=p=0$。此外，经典势小于任意给定数值 $T$ 的区域同样将是紧致的。这些条件意味着，在圆周上量子化时，任意给定能量的量子态在大场值处指数衰减，从而谱是离散的。这样的离散谱随参数连续变化——尽管可能以难以计算的方式变化。

% [Tutor Note: “无 $\sigma$ 场”情形给出了一个可操作的判据：若经典势的零点集紧致，且次能级集 $\{V<T\}$ 对每个 $T$ 都紧致，则圆周量子化所得哈密顿量具有紧预解式，谱离散、波函数在大场处指数衰减。反过来，只要存在一段势为常数的非紧方向，靶空间就不再有效紧致。$r\neq0$ 时 $\sigma$ 正是这样的方向：$s_i=p=0$ 处经典势取值 $D^2/2e^2=e^2r^2/2$，与 $\sigma$ 无关，故 $|\sigma|\to\infty$ 是平坦的；到 $r=0$ 时该常数值本身趋于零，非紧方向与零点集重合，这才是“$r=0$ 处表观奇点”的真正来源。]
% [Tutor Note: 单个能级对参数的连续依赖来自哈密顿量的紧预解式与参数解析性（Kato 摄动理论），但它并不排除能级交叉；受保护的是 $\operatorname{Tr}(-1)^F$ 这类指标——在 $H>0$ 的态成对抵消后它只由 $H=0$ 的态计数给出。Witten 特意加注 "perhaps in a difficult-to-calculate fashion"，提醒的是：谱的连续“存在”与谱的具体数值可计算是两回事，本节只用到前者。]

% ==== 原文 PDF 第 21 页（印刷页 21）====
\label{op:21}
$\sigma$ 场即使在 $r\neq0$ 时也改变了这一图像。虽然对任意 $r$，经典真空的空间都是紧致的（这一点上面已经确定），但经典势小于某个值 $T$ 的区域只在 $T$ 足够小时才是紧致的。事实上，当 $s_i=p=0$ 时，经典势取常数值 $D^2/2e^2=e^2r^2/2$，与 $\sigma$ 无关。这意味着只有谱的低能部分——即能量低到使无穷远区域不可及的那部分——才是离散的。若理论在半径为 $R$ 的圆周上量子化，则在能量高于某个临界值 $T$ 时，理论具有连续谱，它来自支撑在 $s_i=p=0$、$|\sigma|\gg0$ 附近的波函数。半经典地，$T$ 的临界值看来是
\begin{equation}\tag{3.8}
T_{cr}=\frac{e^2r^2}{2}\cdot2\pi R.
\end{equation}

% [Tutor Note: (3.8) 中的 $e^2r^2/2$ 来自消去 $D$：由 FI 项 $-rD$ 与 $D^2/2e^2$ 的 $D$ 运动方程得 $D=-e^2\bigl(\sum_iQ_i|\phi_i|^2-r\bigr)$（同 (3.1)），在 $s_i=p=0$ 处 $D=e^2r$，故 $U=D^2/2e^2=e^2r^2/2$。$T_{cr}=U\cdot2\pi R$ 的含义是：在半径 $R$ 的圆周上量子化时，波函数要扩展到 $|\sigma|\to\infty$ 的区域须付出这一常数势能密度乘以圆周长度的能量代价；这是纯半经典估计，下一段起将检验它对量子修正的稳定性。]

在继续之前，先评估这一估计的可靠性。$T$ 的非零值来自 $s_i,p$ 接近零、$|\sigma|\gg0$ 的区域中 $D$ 场的期望值。在该区域中，包括 $s_i$、$p$ 及其超伴子在内的手征超场具有正比于 $|\sigma|$ 的质量；规范多重态（$\sigma$、$v$ 与 $\lambda$）是无质量的。若简单地把 $s_i$ 与 $p$ 取为零，规范多重态就由自由场论描述，并且（除了我们稍后会涉及的关于 $\theta$ 角的微妙之处外）对经典值 (3.8) 根本没有任何量子修正。朴素地看，积掉质量正比于 $|\sigma|$ 的场只会给出正比于 $|\sigma|$ 的负幂次的修正，不应改变这一结论。只要紫外行为足够好，情况就会如此。

为看清实际发生的情形，考虑一个含电荷为 $Q_i$ 的手征超场 $\Phi_i$ 的一般模型。对 $-\frac{D}{e^2}=\sum_iQ_i|\phi_i|^2-r$ 的真空期望值的单圈修正将是
\begin{equation}\tag{3.9}
\delta\left\langle-\frac{D}{e^2}\right\rangle=\sum_iQ_i\int\frac{\mathrm{d}^2k}{(2\pi)^2}\frac{1}{k^2+2|\sigma|^2+\ldots}.
\end{equation}
（其中 $\ldots$ 是质量中与 $\sigma$ 无关的贡献。）该式对大 $\sigma$ 消失的条件是
% [跨页接续: 原文第21页/第22页]
% ==== 原文 PDF 第 22 页（印刷页 22）====
\label{op:22}
\begin{equation}\tag{3.10}
\sum_iQ_i=0.
\end{equation}
若不满足 (3.10)，则 (3.9) 发散。

% [Tutor Note: 分母中的 $m^2=2|\sigma|^2$ 可由 (3.23) 的势能项 $2|\sigma|^2\sum_i|s_i|^2$ 直接读出（Wess--Zumino 规范下，荷为 $Q_i$ 的手征多重态质量平方为 $2|\sigma|^2$，因子 2 依赖规范多重态的归一化约定，核对符号时应与本节 $L_{gauge}$ 的约定一致）。$\sum_iQ_i=0$ 不只是技术条件：它正是 (2.34) 那种无反常 $R$ 对称性条件。注意不要把它和下面的 $\mathbf{CP}^{n-1}$ 模型混淆：§3.1 的 GLSM（$n$ 个荷 $1$ 的 $S_i$ 加荷 $-n$ 的 $P$）满足 $\sum_iQ_i=0$；而本页末讨论的 $\mathbf{CP}^{n-1}$ 模型只有 $n$ 个荷 $1$ 的场，$\sum_iQ_i=n\neq0$，正属于下一小节“$\sum_iQ_i\neq0$”的情形。$\sum_iQ_i=0$ 成立时单圈修正对大 $\sigma$ 消失，$r$ 不跑动。]

为重整化理论，我们在比如 $|\sigma|=\mu$ 处减去 $\delta D$ 的值。这一减除被解释为参数 $r$ 的重定义。减除之后，对大 $\sigma$ 的行为有
\begin{equation}\tag{3.11}
\begin{aligned}
\delta\left\langle-\frac{D}{e^2}\right\rangle
&=\sum_iQ_i\int\frac{\mathrm{d}^2k}{(2\pi)^2}\left(\frac{1}{k^2+2|\sigma|^2}-\frac{1}{k^2+2\mu^2}\right)\\
&=\sum_iQ_i\cdot\left(\frac{\ln(\mu/|\sigma|)}{2\pi}\right).
\end{aligned}
\end{equation}
该公式在 $|\sigma|\to\infty$ 时发散。（该公式在 $\sigma$ 接近 0 时也发散，但这一效应是误导性的；因为当 $\sigma$ 接近 0 时，我们据以展开的点 $s_i=p=0$ 是不稳定的，近似已不适用。）这一发散意味着：若 $\sum_iQ_i\neq0$，则在这个场空间区域中量子修正是大的，关键性质无法从经典拉格朗日量读出。对 $D$ 的修正可以解释为 $r$ 的一个与 $\sigma$ 有关的有效值：
\begin{equation}\tag{3.12}
r_{eff}\sim r+\frac{1}{2\pi}\sum_iQ_i\cdot\ln(|\sigma|/\mu).
\end{equation}

% [Tutor Note: (3.11) 可直接算出。角向积分后 $\int\frac{\mathrm{d}^2k}{(2\pi)^2}=\frac{1}{2\pi}\int k\,\mathrm{d}k$，而 $\int k\,\mathrm{d}k\left(\frac{1}{k^2+a^2}-\frac{1}{k^2+b^2}\right)=\frac{1}{2}\ln\frac{k^2+a^2}{k^2+b^2}\Big|_0^\infty=\ln(b/a)$，取 $a=|\sigma|$、$b=\mu$ 即得 $\frac{1}{2\pi}\ln(\mu/|\sigma|)$。$\mu$ 的依赖可被 $r$ 的重定义吸收（与四维的量纲嬗变同类），故只有 $\ln|\sigma|$ 的系数是物理的。]

人们可能会问：什么样的超空间相互作用会把这一依赖 $\sigma$ 的有效 $r$ 值作为其后果之一？回顾式 (2.33)，答案很清楚。我们已经生成了一个扭手征超势，其形式恰为 (2.31) 中所考虑过的，其中
\begin{equation}\tag{3.13}
4\pi p=\sum_iQ_i
\end{equation}
——这与我们先前的普遍观察即 $4\pi p$ 必须为整数一致。（在 $\mathbf{CP}^{n-1}$ 模型的情形——我们稍后会专门讨论它——这一效应最先
% [跨页接续: 原文第22页/第23页]
% ==== 原文 PDF 第 23 页（印刷页 23）====
\label{op:23}
由 D'Adda 等人 [21] 从另一观点算出。该普遍公式也已被猜测 [22]。）作者把如下验证留给好奇的读者：通过在大 $\sigma$ 区域的进一步计算，来检验这一扭手征超势同样预言的、依赖 $\sigma$ 的有效 $\theta$ 值。

% [Tutor Note: 与 (2.33) 对照即可定出 $p$：由 (2.33) $r_{eff}=r+2p\ln|\sigma|$，而 (3.12) 给出 $r_{eff}=r+\frac{1}{2\pi}\sum_iQ_i\ln|\sigma|$，故 $2p=\frac{1}{2\pi}\sum_iQ_i$，即 (3.13) $4\pi p=\sum_iQ_i$。这一整数性来自 $\ln\sigma$ 只在模 $2\pi i\mathbb{Z}$ 下有定义，而 $\int_\Sigma\mathrm{d}^2y\,v_{01}\in2\pi\mathbb{Z}$（见 (2.32) 前后的讨论）。]

在讨论这一计算的后果之前，先评估它的相关性。我们刚才算出的这个特定量子修正是否是唯一重要的修正？事实上，对大 $\sigma$，该理论是无质量自由场的理论，这些场通过超可重整的相互作用与质量正比于 $\sigma$ 的场弱耦合。除非存在紫外发散，微扰论对大 $\sigma$ 将一致成立（而且实际上越来越好）。（如我们刚才所见，这类发散经重整化后会转化为随 $\sigma$ 增大而增长的效果。）由于这些模型的超可重整性以及超势的非重整化定理，唯一发散的量子修正就是我们刚刚遇到的 $r$ 的单圈重整化。（在四维可重整超对称理论中一个大致类似的事实是：唯一的二次发散是对 $r$ 的单圈重整化，它同样正比于 $\sum_iQ_i$。）这就是为什么这个修正才是重要的量子修正。

\subsubsection*{$\sum_iQ_i=0$ 的情形}

现在讨论上述计算的后果。先设 $\sum_iQ_i=0$。在这种情形下，对大 $\sigma$ 我们可以干脆扔掉质量量级为 $|\sigma|$ 的手征超场。剩下还须看无质量规范多重态的有效理论；该多重态的有效拉格朗日量为 $L_{gauge}+L_{D,\theta}$。这个有效理论是自由理论，而超对称保证 $\sigma$ 与 $\lambda$ 的零点涨落对基态能量的通常量子修正是抵消的。但对规范场 $v$，会出现一个重要的微妙之处。$v$ 的拉格朗日量为
\begin{equation}\tag{3.14}
L=\int_\Sigma\mathrm{d}^2y\left(\frac{1}{2e^2}v_{01}^2+\frac{\theta}{2\pi}v_{01}\right).
\end{equation}
二维阿贝尔规范理论中 $\theta$ 项的意义
% [跨页接续: 原文第23页/第24页]
% ==== 原文 PDF 第 24 页（印刷页 24）====
\label{op:24}
早已被解释 [23]。它在真空中诱导出恒定的电场，其值等于强度为 $\theta/2\pi$ 的电荷所产生的电场。这一恒定电场对真空能量有贡献。该贡献的值正是人们简单地对 $v_{01}$ 极小化拉格朗日量时所能猜到的值，即 $(e^2/2)\cdot(\theta/2\pi)^2$。更准确地说，只要 $|\theta|\le\pi$ 便如此。否则，通过一个涉及对产生的过程，对能量的贡献可以降低到 $(e^2/2)\cdot((\theta/2\pi)-n)^2$ 在 $n\in\mathbb{Z}$ 上的最小值。因此 $\theta$ 项对真空能量密度的贡献一般为
\begin{equation}\tag{3.15}
\frac{e^2}{2}\left(\frac{\tilde\theta}{2\pi}\right)^2
\end{equation}
其中 $|\tilde\theta|\le\pi$ 且 $\tilde\theta-\theta\in2\pi\mathbb{Z}$。特别地，真空能量密度是 $\theta$ 的连续、周期但不光滑的函数，周期为 $2\pi$。

% [Tutor Note: 二维阿贝尔规范场没有传播自由度，$v_{01}$ 只是一个可积掉的常量：$v_{01}$ 的极值条件给出 $v_{01}=-e^2\theta/2\pi$，代回即得能量密度 $(e^2/2)(\theta/2\pi)^2$。$|\tilde\theta|\le\pi$ 来自真空分支：$\theta\to\theta-2\pi n$ 相当于换一个分支，而对产生（二维瞬子）可使能量弛豫到最近的极小，故能量是 $2\pi$ 周期、连续但不光滑的分段二次函数——与周期势中粒子的能带图像相类比。]

我们由此可以确定 (3.8) 的精确的、经量子修正的类比：大 $\sigma$ 处量子态的最小能量密度为
\begin{equation}\tag{3.16}
U=\frac{e^2}{2}\left(r^2+\left(\frac{\tilde\theta}{2\pi}\right)^2\right).
\end{equation}
在周长为 $2\pi R$ 的圆周上量子化规范多重态所得的量子态的最小能量为
\begin{equation}\tag{3.17}
T'_{cr}=U\cdot2\pi R.
\end{equation}
量子能谱在低于 $T'_{cr}$ 的能量上是离散且连续变化的，因为在这种能量下靶空间是有效紧致的，大 $\sigma$ 区域不可及。只要 $r$ 或 $\theta$ 之一非零，$T'_{cr}$ 就非零，从而存在一段具有这种离散谱的能量范围。

沿着类似思路的另一个论断在我们的一些应用中很重要。可以不去考虑能量 $H$，而考虑 $H+P$，其中 $P$ 为
% [跨页接续: 原文第24页/第25页]
% ==== 原文 PDF 第 25 页（印刷页 25）====
\label{op:25}
动量。由于真空有 $P=0$，而粒子激发对 $H+P$ 的贡献非负，结论相同：当 $H+P<T'_{cr}$ 时谱是离散的。

在我们的所有应用中，当我们试图真正从 Landau--Ginzburg 与 Calabi--Yau 模型的关系中学到些什么时，都会使用某种推广的拓扑推理。在最简单的情形，我们考虑其中一个超荷的指标（等于 $\operatorname{Tr}(-1)^F$，或下文所解释的椭圆亏格（elliptic genus））。已知在某正常数以下 $H$（或 $H+P$）的谱是离散且连续变化的，这些指标便可通过计数 $H=0$（或 $H+P=0$）的态来计算，特别地是良定义且在形变下不变的。因此，除 $r=\theta=0$ 之外，这些量与 $r$、$\theta$ 无关，我们可以取 $\theta\neq0$ 而由 $r>0$ 安全地插值到 $r<0$。

我们还将考虑更精细的量，它们涉及 $H=0$（或 $H+P=0$）的态之间的耦合。其中一些量（例如与 Yukawa 耦合（Yukawa coupling）有关的量）被预言为与 $r$、$\theta$ 无关，另一些则被预言为随 $t$ 全纯地变化。这些预言依赖于最终立足于场空间中分部积分的论证（例如见 [24]）。若路径积分的被积函数在 $\sigma\to\infty$ 时行为很坏，这些论证可能失效。然而，只要 $T'_{cr}\neq0$，$H=0$（或 $H+P=0$）的态的波泛函在 $\sigma\to\infty$ 时便指数衰减。于是这类态的耦合的路径积分表示高度收敛，场空间中的分部积分是有依据的。

总之，对有限体积中的拓扑量而言，唯一真正的奇点位于 $r=\theta=0$。取 $\theta\neq0$，我们便可以从 Calabi--Yau 愉快地延拓到 Landau--Ginzburg。

\subsubsection*{若 $\sum_iQ_i\neq0$}

现在来分析 $\sum_iQ_i\neq0$ 时会发生什么。

在我们的经典拉格朗日量中，规范多重态 $\Sigma$ 的扭超势
% [跨页接续: 原文第25页/第26页]
% ==== 原文 PDF 第 26 页（印刷页 26）====
\label{op:26}
只是一个线性函数 $\widetilde W(\Sigma)=it\Sigma/2\sqrt2$。人们可能会猜想：正如通常的 $W$ 导致玻色势 $\sum_i|\partial_iW|^2$，扭超势也会导致势能
\begin{equation}\tag{3.18}
\left|\frac{\partial\widetilde W(\sigma)}{\partial\sigma}\right|^2.
\end{equation}
这几乎是对的，只差一个来自我们对规范多重态所用归一化的因子。事实上，(3.16) 中的势能是 $U=4e^2|\widetilde W'_Q|^2$，其中 $\widetilde W_Q(x)=i\hat tx/2\sqrt2$；这里 $\hat t=t+n$，$n\in\mathbb{Z}$ 选得使 $U$ 极小。注意，这一（近乎）$|\widetilde W'|^2$ 的项出现在能量中的机制相当微妙：它一部分来自积掉 $D$ 辅助场，一部分来自 $\theta$ 角的效果。

% [Tutor Note: Witten 用 "one might suspect" 一带而过的正是这一步：把扭手征超势的 $F$ 项与规范多重态的 $D$ 项一并消去。由 (2.30)，$\Delta L$ 含 $\sqrt2\widetilde W'(\sigma)(D-iv_{01})$，与 $D^2/2e^2$ 及 $\theta v_{01}/2\pi$ 联立、积掉 $D$ 与 $v_{01}$ 后即得 $|\widetilde W'|^2$ 型的势能；下句给出的归一化是 $U=4e^2|\widetilde W'|^2$，而不是朴素的 $e^2|\widetilde W'|^2$。用 $\widetilde W_Q=i\hat tx/2\sqrt2$ 检验：$U=4e^2|\hat t|^2/8=(e^2/2)|\hat t|^2$，正是 (3.16) 的 $\frac{e^2}{2}(r^2+(\tilde\theta/2\pi)^2)$。漏掉归一化因子会使有效势的整体系数出错。]

公式 (3.18) 出于同样的理由（除同样的修正外）仍然成立，可用来计算由更一般的扭超势
\begin{equation}\tag{3.19}
\widetilde W=\frac{1}{\sqrt2}\left(\frac{it\Sigma}{2}-\frac{\sum_iQ_i}{2\pi}\Sigma\ln(\Sigma/\mu)\right)
\end{equation}
所诱导的势能，后者已包含量子修正。
% [原文疑误: (3.19) 中 $\ln$ 项的系数 $\sum_iQ_i/2\pi$ 与 (3.13) 及 §2 不自洽。按 (2.31) 之前给出的 $\widetilde W(x)=-x\ln x\cdot p/\sqrt2$，(3.19) 给出 $p=\sum_iQ_i/2\pi$，即 $4\pi p=2\sum_iQ_i$；而 (3.13)（由 (2.33) 与 (3.12) 比较而得）给出 $4\pi p=\sum_iQ_i$。若把 (3.19) 的分母 $2\pi$ 改为 $4\pi$，则 $p=\sum_iQ_i/4\pi$，与 (3.13) 及 (2.32) 的整数条件一致；并且代入 Witten 自述的 $U=4e^2|\widetilde W'|^2$ 恰得 (3.20)（其分母确为 $2\pi$），而按印出的 $2\pi$ 则得 $U\propto|\mathrm{i}\hat t-2X|^2$（$X\equiv\frac{1}{2\pi}\sum_iQ_i(\ln(\sigma/\mu)+1)$），其极值条件与 (3.21)、(3.22) 不符。故疑 (3.19) 印误，建议分母改为 $4\pi$；此处照印出的 $2\pi$ 抄录。]
毕竟，对大 $\sigma$，我们只是有一个 $r$ 与 $\theta$ 的有效值带有弱 $\sigma$ 依赖的理论，而我们可以像这些值为严格常数时那样计算能量对 $r$ 与 $\theta$ 的依赖。因此 $\sigma$ 的有效势正是
\begin{equation}\tag{3.20}
U(\sigma)=\frac{e^2}{2}\left|i\hat t-\sum_i\frac{Q_i}{2\pi}\left(\ln(\sigma/\mu)+1\right)\right|^2.
\end{equation}
同样 $\hat t=t+n$，其中 $n\in\mathbb{Z}$ 须选得使 $U(\sigma)$ 极小。因此 $U(\sigma)$ 是 $\sigma$ 与 $t$ 的连续但不光滑的函数。

现在可以看出，当 $\sum_iQ_i\neq0$ 时，关于 $\sigma$ 模型与 Landau--Ginzburg 模型对应关系的分析必须修改，因为对真空结构的半经典判定不再适用。对 $\sum_iQ_i>0$，
% [跨页接续: 原文第26页/第27页]
% ==== 原文 PDF 第 27 页（印刷页 27）====
\label{op:27}
当 $r\ll0$ 时存在半经典推理未曾预言的新真空态；而对 $\sum_iQ_i<0$，这类新态出现在 $r\gg0$ 时。这些新的基态由
\begin{equation}\tag{3.21}
i\hat t-\frac{\sum_iQ_i}{2\pi}\left(\ln(\sigma/\mu)+1\right)=0
\end{equation}
确定，且恰有 $|\sum_iQ_i|$ 个。新的解显式为
\begin{equation}\tag{3.22}
\sigma=\frac{\mu}{e}\exp\left(\frac{2\pi i\hat t}{\sum_iQ_i}\right);
\end{equation}
共有 $|\sum_iQ_i|$ 个，这是由于 $\hat t\to\hat t+n$ 的自由度。当然，(3.21) 中只有出现在大 $\sigma$ 处、即近似成立的那些解才是相关的；这就是为什么需要 $r$ 很大且符号合适。

(3.20) 在无穷远处像 $|\ln\sigma|^2$ 那样增长。这保证在有限世界面体积下，只要 $\sum_iQ_i\neq0$，模型在任何 $r$ 或 $\theta$ 处都没有奇点。\footnote{这与如下事实相符：对于 §3.3 中引入的 A 模型或半扭模型（half-twisted model），当 $\sum_iQ_i\neq0$ 时瞬子求和总是有限的，因为反常的 $R$ 对称性保证任何给定振幅只接收来自有限多个瞬子数取值的贡献。当 $\sum_iQ_i=0$ 时，瞬子求和是无限和，可以带有奇点；我们将在 §5.5 中看到一个例子。}因此，计入无穷远处的额外态，$\sigma$ 模型/L.G. 对应关系的推广即使在 $\sum_iQ_i\neq0$ 时也成立。

% [Tutor Note: $|\ln\sigma|^2$ 的增长正是本节的关键：它把 $r=0$、$\theta=0$ 处仅存的非紧方向也堵住，因此只要 $\sum_iQ_i\neq0$，有限体积下不存在任何 $r$、$\theta$ 处的奇点——这就是“计入无穷远新态后的 C-Y/L-G 对应”。(3.21) 是复方程，其解 (3.22) 只有在大 $\sigma$ 时近似才有效；由 (3.22) 可见 $|\sigma|$ 随 $\hat t$ 沿适当方向取大值而指数增长，故需要 $r$ 大且符号合适。$|\sum_iQ_i|$ 个解来自 $\hat t\to\hat t+n$：$\exp(2\pi i\hat t/\sum_iQ_i)$ 在 $n\to n+|\sum_iQ_i|$ 下不变，而 (2.32) 保证 $\sum_iQ_i\in\mathbb{Z}$，故恰有 $|\sum_iQ_i|$ 个互不相同的解。]

作为例证，来看这些新态如何在 $\mathbf{CP}^{n-1}$ 模型中进入。考虑规范群为 $U(1)$、含 $n$ 个电荷均为 1 的手征超场 $S_i$ 的模型。规范不变性要求超势为 $W=0$。经典势简单地是
\begin{equation}\tag{3.23}
U(s_i,\sigma)=\frac{e^2}{2}\left(\sum_i\bar s_is_i-r\right)^2+2|\sigma|^2\sum_i|s_i|^2.
\end{equation}
对 $r\gg0$，经典真空必须有 $\sigma=0$、$\sum_i|s_i|^2=r$。这样的真空空间在模去规范变换后，是 $\mathbf{CP}^{n-1}$ 的一份拷贝，其凯勒类正比于 $r$，而低能下的模型是以此为靶空间的 $\sigma$ 模型。

% [Tutor Note: (3.23) 就是消去 $D$ 后的经典标量势：由 (3.1) 型公式（此处 $n$ 个场的荷为 $1$、$p=0$）$D=-e^2\bigl(\sum_i|s_i|^2-r\bigr)$ 给出第一项 $D^2/2e^2$；第二项 $2|\sigma|^2\sum_i|s_i|^2$ 来自 (2.19) 中 $\sigma$ 给出的手征多重态质量项，也正是 (3.9) 中 $m^2=2|\sigma|^2$ 的来源。对 $r\gg0$，真空要求 $\sigma=0$ 且 $\sum_i|s_i|^2=r$，故经典真空（模去规范变换后）紧致，这正是本段据以断定低能有效理论为 $\mathbf{CP}^{n-1}$ $\sigma$ 模型的依据。]

% ==== 原文 PDF 第 28 页（印刷页 28）====
\label{op:28}
另一方面，对 $r\ll0$，超对称看来是自发破缺的。这带来一个问题。对 $r\gg0$，理论的超对称指标 $\operatorname{Tr}(-1)^F$\footnote{其定义众所周知，并在下文给出。}等于 $n$（即 $\mathbf{CP}^{n-1}$ 的欧拉示性数），而超对称自发破缺将意味着该指标在 $r\ll0$ 时为 0。鉴于势在无穷远处的 $|\ln\sigma|^2$ 行为，真空态不可能流入或流出无穷远，因而指标应当是常数。为解决这一问题，只需注意到 $\sum_iQ_i=n$。于是所预期的 $n$ 个零能态可通过令 $s_i=0$ 并把 $\sigma$ 取为 (3.21) 的一个解而找到。

$\mathbf{CP}^{n-1}$ $\sigma$ 模型向负 $r$ 的另一种延拓最近被讨论过 [25]。

\subsubsection*{寓意再探}

我们的论证只对粗糙到即使在有限世界面体积中也可见的奇点敏感。无限世界面体积下的相结构如何？Calabi--Yau 与 Landau--Ginzburg 是否被 $r=0$ 处或附近的真正相变所分隔？

没有理由认为这个问题的答案必须是“普适的”，即与在不一定共形不变的多参数理论空间中从 Calabi--Yau 插值到 Landau--Ginzburg 时沿哪条路径无关。沿某条合适的路径，完全可能存在一个刻画分明的相变；而沿另一条路径则可能没有。这看来相当合理。

情形于是将类似于流体的气“相”与液“相”之间的关系。我们将有一个系统，它在一种极限下被很好地描述为 Calabi--Yau，在另一种极限下被很好地描述为 Landau--Ginzburg。但这将是同一系统的两个不同极限，二者之间没有序参量或不可避免的奇点将它们分开；于是人们可以说（按惯例）Landau--Ginzburg 与 Calabi--Yau 是“等价的”。
